\documentclass{beamer}

\usepackage{setspace}
\usepackage{beamerthemesplit}
\usepackage{graphicx}
\usepackage{epstopdf}
\usepackage[latin1]{inputenc}
\usepackage{times}
\usepackage[T1]{fontenc}
\usepackage{color}
\usepackage{colortbl}
\usepackage{multimedia,xmpmulti}
\usepackage{multirow,multicol}
\usepackage[absolute,overlay]{textpos}
\usepackage{tabularx}
\usepackage[noend]{algorithm,algpseudocode}

% Macros (minimal subset -- tensor_header not needed)
\newcommand{\M}[1]{\mathbf{#1}}
\newcommand{\lt}{\left}
\newcommand{\rt}{\right}

%% Tikz
\usepackage{tikz}
\usetikzlibrary{3d}
\usetikzlibrary{patterns}
\usetikzlibrary{calc}
\usetikzlibrary{arrows}

\usepackage{pgfplots}
\usepackage{pgfplotstable}
\usepackage{etoolbox}
\pgfplotsset{compat=1.18}
\usetikzlibrary{decorations.pathreplacing}

\graphicspath{{fig/}}

\mode<presentation>

\definecolor{wfugold}{rgb}{0.6196078,0.494117647,0.21960784}
\usetheme{Warsaw}
\usecolortheme[named=wfugold]{structure}
\setbeamertemplate{navigation symbols}{}
\setbeamertemplate{footline}
{
  \hbox{
  \begin{beamercolorbox}[wd=.33\paperwidth,ht=2.25ex,dp=1ex,left]{author in head/foot}%
    \usebeamerfont{author in head/foot}
    DuBose Tuller
  \end{beamercolorbox}%
  \begin{beamercolorbox}[wd=.34\paperwidth,ht=2.25ex,dp=1ex,center]{title in head/foot}%
    \usebeamerfont{title in head/foot}

  \end{beamercolorbox}%
  \begin{beamercolorbox}[wd=.33\paperwidth,ht=2.25ex,dp=1ex,right]{date in head/foot}%
    \usebeamerfont{date in head/foot}
    \insertframenumber{} \hspace*{2ex}
  \end{beamercolorbox}}%
}

\title{Parallelizing Batched Tensor-Times-Vector on GPUs}
\author{DuBose Tuller}
\institute{CSC 347/647 --- Spring 2026}
\date{April 28, 2026}

\begin{document}

% ============================================================
\begin{frame}
  \titlepage
\end{frame}

% ============================================================
\begin{frame}{Motivation}

\begin{itemize}
  \item Personal research: applying Newton's method to compute the CP decomposition of tensors
  \item The \textbf{batched tensor-times-vector} operation is one of the most expensive in that algorithm
  \item Simple arithmetic: the problem is scale
\end{itemize}

\end{frame}

% ============================================================
\begin{frame}{What is Batched TTV?}

Given a 3rd-order tensor $\mathcal{X} \in \mathbb{R}^{I \times J \times K}$ stored in column-major (natural) order:
\[
  \mathcal{X}(i,j,k) \;=\; \texttt{X[i + j*I + k*I*J]}
\]

\textbf{Tensor-times-vector (mode $m$):}
contract one mode of $\mathcal{X}$ with a vector $\M{u}$, producing a matrix.

\medskip

This operation is equivalent to multiple independent matrix-vector products.

\medskip

\textbf{Example --- contract mode 1, batch over mode 3:}
\[
  Y(j,k) = \sum_{i=1}^{I} \mathcal{X}(i,j,k) \cdot U(i,k)
\]

There are \textbf{six} $(m,n)$ pairings for a 3rd-order tensor.

\end{frame}

% ============================================================
\begin{frame}[fragile]{Serial Implementation}

Implemented here naively as a triple nested loop: two indexings and a dot product.

\medskip

\begin{columns}[T]
\begin{column}{0.60\textwidth}
\begin{algorithmic}[1]
\small
\For{$k = 0 \ldots K{-}1$}
  \For{$j = 0 \ldots J{-}1$}
    \State $s \gets 0$
    \For{$i = 0 \ldots I{-}1$}
      \State $s \mathrel{+}= X[i{+}jI{+}kIJ] \cdot U[i{+}kI]$
    \EndFor
    \State $Y[j{+}kJ] \gets s$
  \EndFor
\EndFor
\end{algorithmic}
\vspace{4pt}
\centering\small Case $(m{=}1, n{=}3)$
\end{column}
\begin{column}{0.45\textwidth}
\hspace{-10pt}
\begin{itemize}\small
  \item All six cases have the same structure
  \item Only the index formulas change
  \item $O(IJK)$ work for each case
\end{itemize}
\end{column}
\end{columns}

\end{frame}

% ============================================================
\begin{frame}{Memory Access Patterns Matter}

$\mathcal{X}$ is column-major: \texttt{X[i + j*I + k*I*J]}.

The \textbf{inner loop variable} determines the stride through memory:

\medskip

\begin{center}
\begin{tabular}{cccc}
\hline
\textbf{Case} & \textbf{Inner var} & \textbf{Stride} & \textbf{Time $n=1000$} \\
\hline
$m{=}1, n{=}3$ & $i$ & 1 & 3.3s \\
$m{=}1, n{=}2$ & $i$ & 1 & 3.4s \\
$m{=}2, n{=}3$ & $j$ & $I$ & 3.4s \\
$m{=}3, n{=}2$ & $k$ & $IJ$ & 3.5s \\
$m{=}2, n{=}1$ & $j$ & $I$ & 6.8s \\
$m{=}3, n{=}1$ & $k$ & $IJ$ & 9.8s \\
\hline
\end{tabular}
\end{center}

\medskip

Same number of FLOPs, but up to $3\times$ difference in wall time due to cache behavior.

Similar pattern to BLAS implementation in MATLAB.

\end{frame}

% ============================================================
\begin{frame}{Parallelization: OpenMP}

The two outer loops are independent (each computes one output element).

\begin{center}
\resizebox{0.95\linewidth}{!}{% ── Figure 2: OMP thread scaling at n=1000 (log-log) ────────
% Requires: \usepackage{pgfplots}  \pgfplotsset{compat=1.18}
\begin{tikzpicture}
\begin{axis}[
    xmode=log, ymode=log,
    log basis x=2, log basis y=2,
    xtick={1,2,4,8,16,32},
    xticklabels={1,2,4,8,16,32},
    xlabel={OMP threads},
    ylabel={Speedup over serial},
    legend pos=north west,
    width=0.9\linewidth,
    height=6.5cm,
    xmin=0.8, xmax=48,
    ymin=0.5, ymax=48,
    ymajorgrids=true, xmajorgrids=true,
]
\addplot[dashed,gray,thick,domain=1:32,samples=2] {x};
\addlegendentry{Ideal}
\addplot[mark=o,blue,thick] table[x=threads,y=m1n3] {fig/fig2.dat};
\addlegendentry{$m{=}1,\,n{=}3$}
\addplot[mark=square*,red,thick] table[x=threads,y=m3n1] {fig/fig2.dat};
\addlegendentry{$m{=}3,\,n{=}1$}
\end{axis}
\end{tikzpicture}
}
\end{center}

\end{frame}

% ============================================================
\begin{frame}{Parallelization: CUDA}

One thread per output element; the inner (contracted) dimension is a serial reduction.

\medskip

% [DIAGRAM: GPU thread/block mapping]
% Suggested diagram: A 2D grid of blocks, each block is T x T threads.
% One axis maps to one free index (e.g., j), the other to the second
% free index (e.g., local_k). The contracted index (i) is the serial
% loop inside each thread.
\begin{center}
  \Large\color{gray} [Diagram: thread $\to$ output element mapping]
\end{center}

\medskip

\begin{itemize}
  \item Block size: $32 \times 32$ threads
  \item Grid: $\lceil D_1 / 32 \rceil \times \lceil D_2 / 32 \rceil$ blocks
  \item Thread $(x,y)$ computes one dot product of length $D_3$
\end{itemize}

\end{frame}

% ============================================================
\begin{frame}{Challenge: Fitting in GPU Memory}

\textbf{First attempt:} generate random data on the GPU using \texttt{curandState\_t}.

\medskip

\begin{itemize}
  \item Each \texttt{curandState\_t} is 48 bytes
  \item For $I{=}J{=}K{=}1000$: need $10^9$ states $= 48$ GB
  \item RTX 2080ti has 8 GB of VRAMx
\end{itemize}

\bigskip

\textbf{Solution:} generate data on the CPU, write to a binary file, and load it.

\begin{itemize}
  \item All three implementations (Serial, OpenMP, CUDA) use the same file
  \item Ensures identical input for fair comparison
\end{itemize}

\end{frame}

% ============================================================
\begin{frame}{Challenge: Tensors Larger than VRAM}

Even with CPU-generated data, $\mathcal{X}$ at $2000^3$ is about 30 GB.

\bigskip

Instead, we tile over the $K$ dimension.

\begin{enumerate}
  \item Query free VRAM with \texttt{cudaMemGetInfo}
  \item Compute how many $K$-slices fit (each slice is $I \times J$ floats)
  \item Loop: copy a chunk to the GPU $\to$ launch kernel $\to$ repeat
  \item One final \texttt{cudaMemcpy} for $Y$ back to host
\end{enumerate}

\medskip

This works because $K$-slices are contiguous in memory (column-major layout).

\end{frame}

% ============================================================
\begin{frame}{Results: Speedup by Case ($N = 1000$)}

\begin{center}
\resizebox{0.85\linewidth}{!}{% ── Figure 1: Speedup over serial at n=1000 ─────────────────
% Requires: \usepackage{pgfplots}  \pgfplotsset{compat=1.18}
\begin{tikzpicture}
\begin{axis}[
    ybar,
    bar width=8pt,
    symbolic x coords={m1n3,m1n2,m2n3,m3n2,m2n1,m3n1},
    xtick=data,
    xticklabels={{$(1,3)$},{$(1,2)$},{$(2,3)$},{$(3,2)$},{$(2,1)$},{$(3,1)$}},
    xlabel={Contraction case},
    ylabel={Speedup over serial},
    legend pos=north east,
    width=0.9\linewidth,
    height=6.5cm,
    ymajorgrids=true,
    ymin=0,
    enlarge x limits=0.15,
]
\addplot[fill=orange!70,draw=orange!90!black] table[x=case,y=omp] {fig/fig1.dat};
\addlegendentry{OMP-32}
\addplot[fill=teal!60,draw=teal!80!black] table[x=case,y=gpu] {fig/fig1.dat};
\addlegendentry{GPU}
\end{axis}
\end{tikzpicture}
}
\end{center}

\begin{itemize}
  \item OMP-32: benefit depends greatly on batch mode.
  \item GPU sidesteps cache issues.
\end{itemize}

\end{frame}

% ============================================================
\begin{frame}{Results: GPU vs.\ CPU Across Problem Sizes}

\begin{center}
\resizebox{0.95\linewidth}{!}{% ── Figure 3: Size scaling for case (m=3,n=1), log-log ──────
% Requires: \usepackage{pgfplots}  \pgfplotsset{compat=1.18}
% Vertical line marks the 8 GB VRAM limit: n^3*4 bytes = 8 GB -> n~1260
\begin{tikzpicture}
\begin{axis}[
    xmode=log, ymode=log,
    xlabel={Tensor mode size $n$ \, (cubic $n\times n\times n$)},
    ylabel={'Kernel' Time (s)},
    legend pos=south east,
    width=0.9\linewidth,
    height=6.5cm,
    xmin=400, xmax=2500,
    xtick={500,1000,1500,2000},
    xticklabels={500,1000,1500,2000},
    ymajorgrids=true, xmajorgrids=true,
    unbounded coords=discard,
]
\addplot[mark=triangle*,green!60!black,thick] table[x=n,y=gpu] {fig/fig3.dat};
\addlegendentry{GPU}
% \addplot[mark=square*,orange!80!black,thick] table[x=n,y=omp] {fig/fig3.dat};
\addlegendentry{OMP-32}
\addplot[mark=o,blue,thick] table[x=n,y=ser] {fig/fig3.dat};
\addlegendentry{Serial}
\end{axis}
\end{tikzpicture}
}
\end{center}

\begin{itemize}
  \item Showing case $(m{=}2, n{=}1)$; GPU timing includes H2D/D2H transfers
  \item GPU uses $K$-chunking beyond the 8\,GB VRAM limit (dashed line)
\end{itemize}

\end{frame}

% ============================================================
\begin{frame}{Conclusions}

\begin{itemize}
  \item Batched TTV on paper is embarrassingly parallel
  \item GPU achieves significant speedup over serial and multi-core OpenMP
  \item Memory access patterns (cache lines / coalescing) affect all implementations
  \item $K$-chunking enables processing tensors larger than GPU memory
\end{itemize}

\bigskip

\textbf{Code:} \texttt{github.com/WFUCSCCho/project-5-DuBose-Tuller}

\end{frame}

\end{document}
